\section{Ficha del proyecto}

\subsection{Nombre del sistema y organizaci\'on cliente}
El sistema a desarrollar se denomina \textbf{Sistema de Gesti\'on para Restaurante con M\'odulo de Inteligencia Artificial} (en adelante, \emph{SGR-IA}). La organizaci\'on cliente es el \textbf{El Balcón Restaurant}, ubicado en la ciudad de Quevedo, provincia de Los R\'ios, Ecuador. El proyecto se desarrolla en el marco de la asignatura Ingenier\'ia de Requerimientos (ISR-401) del programa Ingenier\'ia de Software de la Universidad T\'ecnica Estatal de Quevedo.

\subsection{Problema u oportunidad}
La operaci\'on de un restaurante concentra en un mismo turno la atenci\'on en sala, el registro de pedidos, la preparaci\'on en cocina, el control de insumos, la facturaci\'on y la atenci\'on de pedidos a domicilio. Cuando estos procesos se apoyan en registros manuales y en comunicaci\'on verbal, aparecen demoras en horas pico, susceptibilidad a errores en la transcripci\'on de pedidos y dificultad para consultar el estado real del inventario. La norma ISO/IEC/IEEE 29148:2018 exige identificar a las partes interesadas y establecer el l\'imite del sistema antes de definir requisitos \cite{iso29148}; en consecuencia, esta entrega se limita a acordar la visi\'on, el alcance, el contexto y la planificaci\'on de la elicitaci\'on, sin adelantar soluciones tecnol\'ogicas. La identificaci\'on de partes interesadas y la delimitaci\'on del contexto se tratan en los fundamentos de ingenier\'ia de requisitos \cite{pohl2025}, mientras que la selecci\'on de t\'ecnicas seg\'un la fuente y el tipo de conocimiento se sustenta en el syllabus IREB \cite{ireb2026}.

\emph{Estado de validaci\'on:} el enunciado del problema se formul\'o a partir del dominio de negocio y ser\'a confirmado, corregido o descartado con el responsable de El Balcón Restaurant durante la primera sesi\'on de elicitaci\'on. Ninguna cifra operativa se presenta como dato medido de la organizaci\'on hasta contar con evidencia registrada.

\subsection{Objetivo general}
Caracterizar la operaci\'on real del El Balcón Restaurant y delimitar un sistema de gesti\'on asistido por inteligencia artificial que integre pedidos, cocina, inventario y ventas, con supervisi\'on humana sobre las decisiones cr\'iticas, como base para la especificaci\'on de requisitos de las entregas posteriores del PFC.

\subsection{Objetivos de negocio medibles}
Los siguientes objetivos son metas acordadas para la elicita\'cion y el dise\~no; sus l\'ineas base se obtendr\'an de la observaci\'on y de los registros de la organizaci\'on:

\begin{enumerate}[leftmargin=18pt,itemsep=4pt]
    \item \textbf{Reducir el tiempo de registro de un pedido.} Obtener una reducci\'on de al menos 30\% del tiempo medio de registro respecto de la l\'inea base medida en la sesi\'on de observaci\'on, con el mismo n\'umero de personas en turno.
    \item \textbf{Disminuir los errores de transcripci\'on de pedidos.} Alcanzar una tasa de errores por pedido igual o inferior a 2\% en una muestra de al menos 50 pedidos posteriores a la puesta en marcha del prototipo.
    \item \textbf{Disponer de informaci\'on de inventario actualizada.} Lograr que el 100\% de los insumos cr\'iticos tenga registro de entrada y salida en el sistema, con al menos una actualizaci\'on por turno y una conciliaci\'on semanal contra el conteo f\'isico.
    \item \textbf{Mejorar la trazabilidad de ventas.} Generar un resumen diario de ventas y productos m\'as solicitados sin transcripci\'on manual adicional, disponible para el responsable antes del cierre de caja.
\end{enumerate}

\section{Alcance del sistema}

\subsection{Dentro del alcance}
\begin{itemize}[leftmargin=18pt,itemsep=2pt]
    \item Registro y seguimiento de pedidos de sala y de entrega a domicilio, con estados visibles para sala y cocina.
    \item Gesti\'on de reservas y de solicitudes especiales de los clientes, con registro de la informaci\'on relevante para cocina.
    \item Control de inventario de insumos, entradas, salidas, alertas de reposici\'on y conciliaci\'on semanal.
    \item Registro de compras y de proveedores, con soporte para la comparaci\'on de precios.
    \item Registro de ventas, medios de pago y generaci\'on de res\'umenes diarios y semanales.
    \item M\'odulo de inteligencia artificial para sugerencias de reposici\'on y de productos complementarios, con explicaci\'on de cada sugerencia y confirmaci\'on humana obligatoria.
    \item Gesti\'on de usuarios por rol y respaldo de la informaci\'on operativa.
\end{itemize}

\subsection{Fuera del alcance}
\begin{itemize}[leftmargin=18pt,itemsep=2pt]
    \item Desarrollo de una aplicaci\'on m\'ovil propia para clientes o de un marketplace de delivery externo.
    \item Integraci\'on con pasarelas de pago, bancos o con el sistema de facturaci\'on electr\'onica del SRI.
    \item Automatizaci\'on de la preparaci\'on de alimentos o sustituci\'on del criterio del personal de cocina en decisiones de higiene, alergias o seguridad alimentaria.
    \item Migraci\'on de informaci\'on hist\'orica anterior al inicio del proyecto.
    \item Capacitaci\'on presencial extensiva y soporte t\'ecnico posterior a la entrega acad\'emica.
\end{itemize}

\section{Restricciones principales}
\begin{itemize}[leftmargin=18pt,itemsep=2pt]
    \item \textbf{Restricci\'on legal y \'etica:} el tratamiento de datos del personal y de los clientes se rige por la Ley Org\'anica de Protecci\'on de Datos Personales del Ecuador \cite{lopdp2021} y por el C\'odigo de \'Etica de la ACM \cite{acm2018}. No se publican datos personales, grabaciones ni documentos firmados en el repositorio p\'ublico.
    \item \textbf{Restricci\'on de supervisi\'on:} las recomendaciones del m\'odulo de inteligencia artificial son orientativas; la decisi\'on final permanece en una persona responsable.
    \item \textbf{Restricci\'on econ\'omica y tecnol\'ogica:} el prototipo se desarrolla con recursos acad\'emicos, sin licencias de pago y con dependencia de la conectividad disponible en el local.
    \item \textbf{Restricci\'on de tiempo:} el cronograma de la asignatura finaliza en la semana 17 del SPA 2026--2027; la gesti\'on del trabajo se organiza seg\'un los principios del PMBOK \cite{pmbok2021}.
    \item \textbf{Restricci\'on de Alcance:} la entrega 1 no incluye requisitos detallados ni dise\~no de soluci\'on; esos productos se derivan de la elicitaci\'on planificada en la secci\'on~3.
\end{itemize}

\section{Declaraci\'on de uso de inteligencia artificial}
Se utiliz\'o asistencia de inteligencia artificial como apoyo para la revisi\'on de redacci\'on y la verificaci\'on de la compilaci\'on \LaTeX{}. Herramienta: asistente integrado del entorno OpenCode y revisor autom\'atico independiente; versi\'on: sesi\'on de 19 de noviembre de 2026. Prop\'osito: detectar incoherencias formales y comprobar que el documento regenera el PDF con \texttt{pdflatex}, \texttt{biber} y dos pasadas adicionales de \texttt{pdflatex}. Verificaci\'on: cotejo manual de cada afirmaci\'on contra las fuentes citadas y ejecuci\'on local del procedimiento de compilaci\'on descrito en el archivo \texttt{README.md}. El contenido t\'ecnico y las decisiones de alcance son de autor\'ia del equipo.
